\documentclass{article}
\usepackage{pathfinder-readable}

\title{When reviewers agree without finding the bug}

\begin{document}
\maketitle

\begin{abstract}
This account connects a paper on incentives for AI agents with a paper on reviewer--critic code review. The
thread asked whether incentives, preserved disagreement and checks on code evidence could prevent reviewers from
agreeing on a wrong answer. The peers established limits on what peer scoring can reveal and derived a condition
under which filtering review flags would improve F1. Neither result shows that a new protocol improves code
review in the measured setting. The thread paused because the needed independent judgments, defect labels and
checks of cited code have not been collected.
\end{abstract}

\section{Introduction}

Bergemann, Koh and Morris study how to direct AI agents whose preferences and abilities are unknown \cite{Q}.
Their mechanisms reward truthful reports and obedience. One example scores a report against peer reports. Another
lets a weaker monitor influence a stronger actor. Both need stated assumptions about beliefs, rewards and actions.

Qiu and Gill study Adversarial Review, a code-review procedure in which a critic challenges a review before a
coding agent acts \cite{P}. On a review-only benchmark, the initial procedure scores \(0.457\) in matched-comment
F1, a score combining precision and recall. The paper describes speculative concerns that survive agreement and a
manually verified concern that
disappears when the critic yields. Revised text instructions score \(0.533\) on the same benchmark.

The connection was a question about false agreement. Could incentives explain it and suggest a way to keep useful
disagreement? The peers separated truthful reporting from finding a real defect. The first paper speaks to
reporting under its assumptions. Defect detection needs evidence about the code.

\section{What was done}

Emmy examined peer scoring; Ada examined the weak-monitor example. In the latter, more control can help a biased
monitor and harm the human \cite{Q}. Ada used this as an analogy for the lost concern, not as a code-review
theorem. Both peers checked the assumptions. The code-review paper reports neither conditional peer beliefs nor
enforceable rewards. Its critic sees the reviewer's judgment before responding \cite{P}.

Ada next modelled a gate that filters review flags. She checked the revised prompts: they combine verdict labels,
code citations and flag-preservation rules, so the F1 gain cannot identify which change helped. Emmy then checked
the peer-score incentive and whether an elicited report would say anything about a bug. The peers tested both
directions with examples. A prior-work check found peer-score training, limited ground-truth checks and
executable code-review verification \cite{peerwork,spotaaai,spotarxiv,codecheck}. This account relies only on the
ledger's check of those sources.

\section{Results}

The incentives paper's peer result concerns truthful reporting by types that the model assigns positive
probability \cite{Q}. For an agent's private type \(t\), let \(\beta[t]\) be its conditional probability
distribution over the other agents' types. If two possible types \(t\) and \(\hat t\) give different
distributions, the quadratic peer score gives truthful reporting an expected advantage of
\[
\Lambda\lVert\beta[t]-\beta[\hat t]\rVert_2^2.
\]
Here \(\Lambda\) is the score weight, and the squared norm measures the difference between the two
distributions. The peers checked that, if the smallest such squared difference is \(\delta>0\) and all other
gains from a false report are at most \(B\), then \(\Lambda\delta\geq B\) blocks that false report. The paper's
full construction also needs known conditional beliefs, enforceable rewards and penalties for actions outside
the target. It implements the desired rule for supported types in an equilibrium; it does not establish that
every possible equilibrium has that outcome.

The peers then showed why this reporting result is separate from defect detection. Let \(Y\) be a fair bit
saying whether there is a bug, and let \(Z\) be an independent fair bit with no bearing on the bug. If both
agents' private types equal \(Z\), each type predicts the other's perfectly. The squared difference \(\delta\)
is \(2\), its maximum for two possible peer types, yet the reports give no information about \(Y\). Conversely,
if one agent's type equals \(Y\) and the other's type is an independent fair bit, the first agent knows the bug
state but its beliefs about the peer do not differ across its own types: \(\delta=0\). Emmy constructed these
examples and Ada checked the first. They are examples inside the first paper's type-and-belief model, not
measurements of the reviewers in the second paper. Peer predictability is therefore neither sufficient nor
necessary for information about defects.

Ada's filtering calculation gives a second, narrower result. Let \(\pi\) be the share of cases with a defect,
\(t\) the original true-positive rate and \(f\) the original false-positive rate. Let a gate retain a fraction
\(s\) of true flags and a fraction \(z\) of false flags. With a positive original true-positive count,
substituting these rates into the F1 formula shows that the gate improves F1 exactly when
\[
(1-\pi)f(s-z)>\pi(1-s).
\]
The left side is the benefit from rejecting false flags more often than true ones. The right side is the cost of
losing true flags. A citation request that keeps both kinds at the same rate cannot meet this condition. This
calculation covers filtering existing flags; the review procedure can also create new flags during discussion.

\section{Objections and limits}

The peers did not infer that using the same model for two roles violates the incentive paper's separation
condition. They found no measurement of that condition in the code-review paper. Nor did they treat a requested
code citation as a checked certificate: the revised prompts ask for citations, but do not verify that the cited
code supports the claimed defect. The revised protocol's \(0.533\) F1 belongs to a review-only benchmark. The
paper's separate repository-repair result uses an earlier form of the protocol \cite{P}.

The review benchmark matches generated comments to human comments. That measure can miss a valid finding absent
from the human review \cite{P}. The papers also do not provide independent judgments from reviewer and critic
before dialogue, a measured record of when true dissent turns into false agreement, or blinded decisions on
whether citations support real defects. These gaps prevent a causal claim that preserving dissent or checking
evidence would improve review quality. The peers also left open whether cheap checks can cover enough findings
and whether any review gain would carry over to successful repairs.

\section{Why the thread stopped}

The verifier paused the thread after one round. The counterexamples and F1 calculation were sound, but they
establish limits and a conditional prediction, not an observed improvement to the reviewer--critic procedure.
The smallest restart is a study on held-out pull requests that records reviewer and critic bug judgments
independently before dialogue. It would vary whether initial dissent is preserved and whether cited code
receives an independent check. Blinded adjudication of defects and citation support would then show whether
these changes reduce false agreement, and at what review and audit cost. Peer scoring would require a further
test of its belief and incentive assumptions before it could support a claim about this procedure.

\bibliographystyle{plainurl}
\bibliography{references}

\end{document}
